\documentclass[11pt]{article}

\usepackage[T1]{fontenc}
\usepackage{lmodern}
\usepackage[a4paper,margin=18mm]{geometry}
\usepackage{microtype}
\usepackage[hidelinks]{hyperref}

\input{metadata}

\hypersetup{
  pdftitle={Cover letter for \papertitle},
  pdfauthor={Marcus Gawronsky},
  pdfsubject={Submission to Quantitative Finance}
}

\pagestyle{empty}
\setlength{\parindent}{0pt}
\setlength{\parskip}{4.4pt}
\setlength{\emergencystretch}{1em}
\linespread{0.966}
\hyphenpenalty=10000
\exhyphenpenalty=10000

\begin{document}

\fontsize{9.85pt}{11.6pt}\selectfont

\begin{flushright}
  30 August 2026
\end{flushright}

Dear Editors of \textit{Quantitative Finance},

Please consider our manuscript, ``\textbf{\papertitle},'' for publication in
\textit{Quantitative Finance}.

Financial models of dependence typically begin only after a firm has been
reduced to a point-valued characteristic vector, assigned a position in a
supplied network, or represented through an estimated return covariance.
Yet economically heterogeneous firms do not occupy single positions: their
operations span products, technologies, counterparties, markets, and information states.
Modern language-model representations now make it possible to preserve much
of this heterogeneity empirically, by representing a firm through a
distribution of information rather than a single score or embedding.
This creates a timely second-order question for quantitative finance:
\textit{once rich distribution-valued representations of firms are available,
what can they establish about financial dependence rather than merely prediction?}

Our paper develops one answer at the pairwise level.
We represent firm characteristics as probability laws and use quadratic
optimal transport to compare them.
The Wasserstein construction has an economic interpretation beyond generic
similarity: it measures the least aggregate displacement required to align
two heterogeneous characteristic distributions, treating substitutions
between nearby economic positions as less costly than substitutions between
distant ones.
In latent factor-risk coordinates, quadratic transport has a further special
property: the squared displacement minimized by $W_2$ is exactly the quantity
that polarization converts into an inner product.
Consequently, for given marginal exposure laws, the optimal transport problem
delivers the \textbf{sharp maximum systematic covariance attainable over
their possible joint arrangements}, while reflection gives the corresponding
lower endpoint.

We then ask when observable characteristic geometry can discipline this
latent covariance set.
Under an explicit common randomized bi-Lipschitz information-to-exposure map,
bounded firm-specific slack in risk coordinates, and a maintained return
bridge, observed distributional separation yields a conditional restriction
on the covariance ceiling.
Economically, if two firms remain far apart under every admissible
transmission of their characteristics into systematic exposures, their risks
cannot be made arbitrarily well aligned.
The resulting object is therefore not another covariance estimator or
expected-return signal.
It is an \textbf{external restriction on the set of systematic dependence
configurations consistent with observable economic composition}.

This perspective is potentially useful where conventional dependence
estimation is weakest: new or thinly traded securities, private or
infrequently valued assets, changing investment universes, and other settings
in which rich information about economic composition is available but joint
return histories are short or heavily proxy-dependent.
More generally, the distributional primitive need not be textual: analogous
characteristic laws could describe geographic footprints, property
portfolios, loan books, supplier or customer exposures, product mixes, or holdings.
In such applications, the covariance envelope could complement rather than
replace factor, shrinkage, or stress-testing approaches by supplying
information about admissible dependence that originates outside joint returns.

Our empirical application uses distributions of language-model
representations of financial news, disclosures, and analyst reports for Nasdaq firms.
Greater pairwise Wasserstein separation is associated with weaker return
co-movement after symmetric firm effects; a one-standard-deviation increase
in $W_2$ corresponds to approximately eight fewer correlation points at
representative dyads.
The association retains the same sign when the information geometry is held
fixed and returns are measured in a later window, and it remains positive
across a substantial representation-sensitivity analysis.
We treat this evidence as reduced-form support for the relevance of the
measured geometry, rather than as causal or point-in-time predictive validation.

We believe the manuscript is particularly suited to \textit{Quantitative
Finance} because it connects factor-model covariance, optimal transport,
spatial notions of economic proximity, and modern representation methods to
introduce a new quantitative object: a \textbf{covariance envelope derived
from distribution-valued firm information}.
The contribution is not that language models predict another financial
variable, but that rich representations can, under explicit economic
conditions, restrict the systematic risk configurations that remain feasible.

We confirm that the manuscript is original, has not been published
previously, and is not under consideration for publication elsewhere.

Thank you for considering our manuscript.

\begingroup
\setlength{\parskip}{0pt}
Sincerely,\par
\smallskip
Marcus Gawronsky\par
Department of Finance and Tax\par
University of Cape Town
\endgroup

\end{document}
